\subsection*{Justificacion Fila 8:}

Primero se proyecta la columna $A$ de $R$. Eliminamos los renglones duplicados, en este caso eliminamos un $2$ ya que aparece dos veces.

\begin{center}
$\pi_{A}(R) =$ \quad
\begin{tabular}{|c|}
\hline
\textbf{A} \\ \hline
  1 \\
  2 \\
  3 \\
  9 \\ \hline
\end{tabular}
\end{center}

Después se proyecta la columna $D$ de $S$. De igual forma, eliminamos un $3$ ya que aparece dos veces.

\begin{center}
$\pi_{D}(S) =$ \quad
\begin{tabular}{|c|}
\hline
\textbf{D} \\ \hline
  3 \\
  1 \\
  0 \\
  2 \\ \hline
\end{tabular}
\end{center}

Se renombra la columna $D$ por $A$ para poder aplicar la intersección, ya que se requiere que los atributos tengan el mismo nombre.

\begin{center}
$\rho_{A \leftarrow D}(\pi_{D}(S)) =$ \quad
\begin{tabular}{|c|}
\hline
\textbf{A} \\ \hline
  3 \\
  1 \\
  0 \\
  2 \\ \hline
\end{tabular}
\end{center}

Se aplica la intersección, que conserva únicamente los valores que aparecen en ambas tablas. Por lo que como $9$ solo está en la primera y $0$ solo en la segunda, estos valores los eliminamos.

\begin{center}
$\pi_{A}(R) \cap \rho_{A \leftarrow D}(\pi_{D}(S)) =$ \quad
\begin{tabular}{|c|}
\hline
\textbf{A} \\ \hline
  1 \\
  2 \\
  3 \\ \hline
\end{tabular}
\end{center}

Por lo tanto, la cardinalidad del resultado es \textbf{3}.